\section*{Introduction}

Time-of-flight positron emission tomography (TOF-PET) uses the measured difference in photon arrival times to constrain the annihilation position 
along the line of response. 
Better timing resolution therefore improves spatial localization and image signal-to-noise performance~\cite{surti2015}. 
Machine-learning (ML) methods can use waveform information beyond a single threshold crossing and have been investigated 
for timing estimation and timing correction in PET detectors~\cite{berg2018,onishi2022}.

The objective of this work is to test a physically structured waveform correction. Leading-edge discrimination (LED) provides 
the conventional timing estimate. A shared detector model \(g_\theta\) is applied independently to the two aligned waveforms 
\(s_1\) and \(s_2\), and the pair correction is constrained to be antisymmetric:
\begin{equation}
    \widehat{y}(s_1,s_2)
    =
    g_\theta(s_1)-g_\theta(s_2).
\end{equation}
Thus,
\begin{equation}
    \widehat{y}(s_2,s_1)=-\widehat{y}(s_1,s_2).
\end{equation}

The detector scorer is implemented with locally connected temporal layers. These layers preserve local receptive fields but do not 
share the same weights across temporal positions, which is useful when the waveforms are aligned and absolute time within the waveform 
has physical meaning~\cite{elsayed2020,amit2019}.

This report is organized to desribe the full scientific pipeline: acquisition, conversion, event selection, ML dataset construction, 
model definition, blind evaluation, and coincidence timing resolution (CTR) results.


